\documentclass[12pt]{article}
\usepackage[a4paper, margin=2cm]{geometry}

\usepackage{fontspec}
\usepackage{polyglossia}
\setdefaultlanguage{russian}
\setmainlanguage{russian}
\setmainfont{Times New Roman}

\begin{document}

\begin{center}
\textbf{\LARGE СЛУЖЕБНАЯ ЗАПИСКА}
\end{center}

\vspace{1cm}

% === ТЕКСТ ДОКУМЕНТА ВСТАВЛЯТЬ СЮДА ===
Кому: Ивану Иванову, начальнику отдела по международному экспорту нефти \\
От: Александр Александрович, сотрудник отдела экспорта, Санкт-Петербург \\
Тема: Предложение по корректировке подхода к планированию поставок

Уважаемый Иван Иванович,

В целях повышения точности планирования поставок предлагаю скорректировать подход к формированию плановых показателей. В настоящее время мы используем усреднённые цены, что приводит к расхождениям между плановыми и фактическими значениями. Предлагаю ориентироваться на актуальную рыночную стоимость нефти Brent на дату формирования планов, чтобы повысить точность расчётов и снизить риски.

На текущий момент Brent торгуется примерно 64,82 доллара за баррель (данные Bloomberg Energy). Привязка плановых показателей к текущей рыночной цене Brent позволит более точно отражать стоимость поставок и связанных с ними обязательств, а также снизить риск несоответствий между бюджетом и фактическими расходами.

Реализация предложения может быть осуществлена следующими шагами:
- определить единый источник котировок Brent и закрепить его в регламенте планирования;
- внедрить использование рыночной цены Brent в процессы планирования и закупок;
- обновлять цену ежедневно в условиях динамичного рынка;
- пересмотреть методику бюджетирования и расчета плановых показателей с привязкой к текущей рыночной цене;
- обеспечить согласование изменений с финансовым отделом.

Эффект от внедрения: повышение точности планирования поставок и снижение рисков, связанных с колебаниями цен на нефть.

С уважением, \\
Александр Александрович \\
Сотрудник отдела экспорта, Санкт-Петербург

\vfill

\noindent
Дата: 25.03.2026 \hfill Подпись: \rule{5cm}{0.4pt}

\end{document}